\documentclass[11pt,a4paper]{article}

\usepackage[T1]{fontenc}
\usepackage[utf8]{inputenc}
\usepackage[margin=2.6cm]{geometry}
\usepackage{amsmath,amssymb}
\usepackage{booktabs}
\usepackage{longtable}
\usepackage{listings}
\usepackage{xcolor}
\usepackage{graphicx}
\usepackage[hidelinks]{hyperref}

\definecolor{shell}{gray}{0.96}
\lstdefinestyle{sh}{
  basicstyle=\ttfamily\small, backgroundcolor=\color{shell},
  frame=single, framerule=0pt, breaklines=true, columns=fullflexible,
  xleftmargin=4pt, xrightmargin=4pt, aboveskip=8pt, belowskip=8pt,
}
\lstdefinestyle{formula}{
  basicstyle=\ttfamily\small, breaklines=true, columns=fullflexible,
  xleftmargin=10pt, aboveskip=6pt, belowskip=6pt,
}
\lstset{style=sh}

\newcommand{\tool}{\texttt{tctl\_sat}\xspace}
\usepackage{xspace}

\title{\tool --- User Manual\\[2pt]
  \large A tableau decision procedure for TCTL satisfiability}
\author{}
\date{}

\begin{document}
\maketitle

\noindent
\tool answers one question: given a TCTL formula, \emph{is there a timed
automaton that satisfies it?}  On success it can also hand back such an
automaton as a certificate.  This manual covers the syntax it accepts, how to
run it, and what it produces.

\tableofcontents

\section{Getting started}

\subsection{Building}

\tool needs a C++20 compiler, CMake $\geq 3.20$,
Z3, and \texttt{git}.

\begin{lstlisting}
cmake -S . -B build -DCMAKE_BUILD_TYPE=Release
cmake --build build -j
\end{lstlisting}

\noindent
That is the whole process.  The zone layer (HDBM) is header-only and is cloned
automatically on the first configure if it is not already present.  Z3 is the
one external library and comes from the system: \texttt{brew install z3} on
macOS, \texttt{apt install z3 libz3-dev} on Linux.

\subsection{A first check}

\begin{lstlisting}
$ ./build/tctl_sat 'EF[1,2] p'
1: SATISFIABLE
\end{lstlisting}

\noindent
Each answer is \texttt{SATISFIABLE}, \texttt{UNSATISFIABLE} or
\texttt{TIMEOUT}, prefixed by the input's line number.

\subsection{What the answers mean}

The procedure is \emph{sound} and \emph{terminating} for all of TCTL: a
\texttt{SATISFIABLE} verdict always means a witness automaton exists.  It is
\emph{complete} over the non-punctual fragment --- every interval $[a,b]$ with
$a<b$ --- and also at temporal nesting depth one even with punctual intervals
$[a,a]$.

Outside that, a punctual interval nested inside another timed modality is
accepted and warned about:

\begin{lstlisting}
$ ./build/tctl_sat 'EG[1,5] (EF[2,2] p)'
1: warning: punctual interval [a,a] nested inside another timed modality ...
1: SATISFIABLE
\end{lstlisting}

\noindent
There, \texttt{SATISFIABLE} is still reliable; \texttt{UNSATISFIABLE} may be a
false negative.

\section{Syntax}\label{sec:syntax}

\subsection{Grammar}

\begin{center}
\begin{tabular}{rcl}
$\varphi$ &::=& \texttt{true} $\mid$ \texttt{false} $\mid$ $p$
              $\mid$ \texttt{\{} \textit{arith} \texttt{\}} \\
          &$\mid$& \texttt{!}$\varphi$ $\mid$ $\varphi$ \texttt{\&} $\varphi$
              $\mid$ $\varphi$ \texttt{|} $\varphi$
              $\mid$ $\varphi$ \texttt{->} $\varphi$
              $\mid$ $\varphi$ \texttt{<->} $\varphi$ \\
          &$\mid$& \texttt{EF}$I\,\varphi$ $\mid$ \texttt{AF}$I\,\varphi$
              $\mid$ \texttt{EG}$I\,\varphi$ $\mid$ \texttt{AG}$I\,\varphi$ \\
          &$\mid$& \texttt{E(}$\varphi$ \texttt{U}$I$ $\varphi$\texttt{)}
              $\mid$ \texttt{A(}$\varphi$ \texttt{U}$I$ $\varphi$\texttt{)} \\
          &$\mid$& \texttt{E(}$\varphi$ \texttt{R}$I$ $\varphi$\texttt{)}
              $\mid$ \texttt{A(}$\varphi$ \texttt{R}$I$ $\varphi$\texttt{)} \\[4pt]
$I$       &::=& $\varepsilon$ $\mid$ \texttt{[}$a$\texttt{,}$b$\texttt{]}
              $\mid$ \texttt{(}$a$\texttt{,}$b$\texttt{)}
              $\mid$ \texttt{[}$a$\texttt{,}$b$\texttt{)}
              $\mid$ \texttt{(}$a$\texttt{,}$b$\texttt{]}
              $\mid$ \texttt{[}$a$\texttt{,inf)} \\
\end{tabular}
\end{center}

\noindent
Atoms are identifiers such as \texttt{p}, \texttt{q}, \texttt{p1}.  An omitted
interval means $[0,\infty)$, so \texttt{EF p} and \texttt{EF[0,inf) p} are the
same formula.

\subsection{Operators}

\begin{center}
\begin{tabular}{llll}
\toprule
Written & Means & Timed form & Reading \\
\midrule
\texttt{EF} $\varphi$ & $\exists\Diamond$ & \texttt{EF[1,5] p} & some run reaches $p$ within $[1,5]$ \\
\texttt{AF} $\varphi$ & $\forall\Diamond$ & \texttt{AF[0,3] p} & every run reaches $p$ within $[0,3]$ \\
\texttt{EG} $\varphi$ & $\exists\Box$     & \texttt{EG[1,5] p} & some run keeps $p$ throughout $[1,5]$ \\
\texttt{AG} $\varphi$ & $\forall\Box$     & \texttt{AG[0,2] p} & every run keeps $p$ throughout $[0,2]$ \\
\texttt{E(}$\varphi$ \texttt{U} $\psi$\texttt{)} & $\exists\,\mathbf{U}$ &
  \texttt{E(p U[1,3] q)} & some run holds $p$ until $q$, in $[1,3]$ \\
\texttt{A(}$\varphi$ \texttt{R} $\psi$\texttt{)} & $\forall\,\mathbf{R}$ &
  \texttt{A(p R[2,5] q)} & $\psi$ holds until released by $\varphi$ \\
\bottomrule
\end{tabular}
\end{center}

\paragraph{There is no next operator.}
TCTL has none, and \texttt{EX}/\texttt{AX} are rejected by the parser.  They
exist only as internal markers inside the construction.

\subsection{Intervals}

Bounds are non-negative integers.  Both closed and open endpoints are supported,
and the strictness is respected exactly: \texttt{EF[1,2] p} and
\texttt{EF(1,2) p} are different questions.

A \emph{punctual} interval $[a,a]$ is accepted --- see the completeness note in
Section~\ref{sec:syntax}'s preceding chapter.  An empty interval (a lower bound
above the upper one, or an open singleton) is a parse error.

\subsection{Arithmetic leaves}

Constraints in braces are handed to Z3 rather than to the temporal machinery:

\begin{lstlisting}[style=formula]
{x <= 5} & {x > 10}          % unsatisfiable: no such x
EF ({x <= 10} & p)           % satisfiable
\end{lstlisting}

\noindent
Variables are integer-sorted by default; \texttt{--real-vars x,y} makes the
named ones reals.

\section{Running the tool}

\subsection{Input}

A formula can be given directly, or in a file with one formula per line:

\begin{lstlisting}
./build/tctl_sat 'AG[0,2] p'          # a single formula
./build/tctl_sat properties.txt       # a file of them
\end{lstlisting}

\noindent
In a file, blank lines are skipped and everything after \texttt{\#} is a
comment.  Each formula is answered independently, prefixed by its line number:

\begin{lstlisting}
# properties.txt -- a small specification
EF[1,3] request                  # a request is possible in [1,3]
AG (request -> AF[0,5] grant)    # every request is granted within 5
AG[0,10] !error                  # no error in the first 10 units
\end{lstlisting}

\begin{lstlisting}
$ ./build/tctl_sat properties.txt
1: SATISFIABLE
2: SATISFIABLE
3: SATISFIABLE
\end{lstlisting}

\noindent
Answering each line separately asks whether each property is satisfiable
\emph{on its own}.  To ask whether they are \emph{jointly} satisfiable --- the
usual specification sanity-check --- conjoin them with \texttt{--tot}:

\begin{lstlisting}
$ ./build/tctl_sat --tot properties.txt
1: SATISFIABLE
\end{lstlisting}

\noindent
This is how an inconsistent specification is caught.  Each of the following is
satisfiable alone, but not together:

\begin{lstlisting}
$ printf 'EF[1,2] p\nAG[0,3] !p\n' > bad.txt
$ ./build/tctl_sat bad.txt
1: SATISFIABLE
2: SATISFIABLE
$ ./build/tctl_sat --tot bad.txt
1: UNSATISFIABLE
\end{lstlisting}

\subsection{Options}

\begin{center}
\begin{tabular}{ll}
\toprule
Option & Effect \\
\midrule
\texttt{--tot}, \texttt{--all} & conjoin every input formula and decide the conjunction \\
\texttt{--model} & extract a witness automaton (Section~\ref{sec:witness}) \\
\texttt{--stats} & report the number of tableau nodes explored \\
\texttt{--output <dir>} & directory for extracted witnesses \\
\texttt{--real-vars a,b} & treat the named arithmetic variables as reals \\
\texttt{--selftest} & run the built-in semantic test suite \\
\texttt{--help} & full option list \\
\bottomrule
\end{tabular}
\end{center}

\subsection{Environment flags}

\begin{center}
\begin{tabular}{ll}
\toprule
Variable & Effect \\
\midrule
\texttt{TCTL\_DEBUG=1} & trace every expanded node and each component, to stderr \\
\texttt{TCTL\_LEGACY\_LOOP=1} & use the earlier path-local loop acceptance (comparison only) \\
\bottomrule
\end{tabular}
\end{center}

\section{Examples}

Each example below was run with the version of \tool this manual ships with.

\subsection{Reachability and safety in conflict}

\begin{lstlisting}
$ ./build/tctl_sat 'EF[1,2] q & AG[0,3] !q'
1: UNSATISFIABLE
\end{lstlisting}

\noindent
The first conjunct demands $q$ somewhere in $[1,2]$; the second forbids it
everywhere in $[0,3]$.  Since $[1,2]\subseteq[0,3]$ there is no model.  Widen
the safety window so the two no longer overlap and a model appears:

\begin{lstlisting}
$ ./build/tctl_sat 'EF[2,2] q & AG[0,1] !q'
1: SATISFIABLE
\end{lstlisting}

\subsection{Strictness matters}

\begin{lstlisting}
$ printf 'EF[1,2] p & AG[0,1] !p\nEF(1,2) p & AG[0,1] !p\n' > s.txt
$ ./build/tctl_sat s.txt
1: SATISFIABLE
2: SATISFIABLE
\end{lstlisting}

\noindent
Both hold: $p$ can be placed strictly after $1$.  The zone layer tracks the
open/closed endpoint exactly rather than rounding it to an integer.

\subsection{A leads-to specification}

The classic response property, \texttt{AG(request -> AF[0,5] grant)}, is the
shape most often written for \textsc{Uppaal}.  It is satisfiable on its own, and
stays so together with a bounded-response requirement:

\begin{lstlisting}
$ ./build/tctl_sat --tot leadsto.txt
1: SATISFIABLE
\end{lstlisting}

\noindent
where \texttt{leadsto.txt} contains

\begin{lstlisting}[style=formula]
AG (request -> AF[0,5] grant)
AG[0,20] !error
EF[1,3] request
\end{lstlisting}

\subsection{Nested timed modalities}

Nesting is where the procedure does real work.  The following is satisfiable,
and the node count shows the cost:

\begin{lstlisting}
$ ./build/tctl_sat --stats 'EG[1,5] (EF[2,4] (EG[6,8] p))'
1: SATISFIABLE
  Stats: nodes=261
\end{lstlisting}

\subsection{A contradiction the tool must catch}

\begin{lstlisting}
$ ./build/tctl_sat 'EF p & !EF p'
1: UNSATISFIABLE
\end{lstlisting}

\section{Witness automata}\label{sec:witness}

A \texttt{SATISFIABLE} verdict is a claim that a model exists.  With
\texttt{--model}, \tool hands the model over as well, so the claim can be
inspected or checked by another tool.

\begin{lstlisting}
$ ./build/tctl_sat --model 'EF[1,2] p'
1: SATISFIABLE
  Witness: 5 locations, 6 transitions, 2 clocks
  wrote witness_1.{dot,json,xml}
\end{lstlisting}

\noindent
Three files are written per satisfiable formula: Graphviz \texttt{.dot} for
viewing, \texttt{.json} for further processing, and \texttt{.xml} in
\textsc{Uppaal}'s format.

\subsection{Reading the output}

For \texttt{EF[1,2] p} the extracted automaton is

\begin{lstlisting}
L0 -> L1 [x_1 < 1]
L1 -> L3 [x_1 <= 1 && x_1 >= 1]
L3 -> L4 [x_1 >= 1]
L4 -> L4 [x_1 >= 1]
L0 -> L2 ; L2 -> L2
\end{lstlisting}

\noindent
with \texttt{p} holding at \texttt{L3}.  Read left to right: let the clock run
below $1$, step onto $x_1 = 1$ --- inside the window $[1,2]$ --- where $p$
holds, then continue on a self-loop so that time diverges.  The
\texttt{L0}--\texttt{L2} branch is the other half of the cover rule, where the
obligation is discharged rather than pursued.

The self-loop at the end matters: a witness must let time diverge, or the
run would be Zeno and the semantics would not inspect it.

\subsection{What is guaranteed}

Locations are the marked nodes, labelled by the atoms they assert and carrying
the deadline of any universal until as an invariant.  Transitions are the
retained discrete successors, guarded by the target's zone, resetting a
formula-clock exactly where its modality is launched.  Delay contributes no
transition: time simply advances inside a location.

Every extracted automaton is checked before it is written --- that it is
non-blocking, that no location asserts both a proposition and its negation, and
that guards only read declared clocks.  A failure is reported as
\texttt{[INVALID]} together with the reason.

\subsection{Limits}

Extraction is a property of the component-based acceptance engine, which is the
default, so it needs no extra flag.  It is the \emph{legacy} engine that cannot
do it: \texttt{TCTL\_LEGACY\_LOOP=1} returns a bare verdict and never builds a
marking, so there is nothing to read an automaton off.  The zone recorded for a node is
the $\cong_M$-canonical one, so guards are read off the extrapolated zone rather
than the raw one --- the same zone the procedure itself compares on.


\end{document}
